% ============================================================================
%  Rozdział: Wyniki treningu oraz testów na rzeczywistej kamerze
%  (modele GRU, LSTM, SSM operujące na narożnikach szachownicy)
%  Plik przeznaczony do samodzielnej kompilacji lub włączenia przez \input{}
%  do głównego dokumentu pracy magisterskiej.
% ============================================================================
% (Uwaga) Ten plik jest włączany do main.tex.
% Dlatego nie może zawierać \documentclass, \usepackage ani \begin{document}/\end{document}.

\section{Wyniki treningu oraz testów na rzeczywistej kamerze modeli GRU, LSTM i SSM}
\label{chap:wyniki-gru-lstm-ssm}

\subsection{Metodyka eksperymentu}
\label{sec:wyniki-metodyka}

Trzy modele sekwencyjne opisane w rozdziale ---
\texttt{FisheyeCornerGRUSequenceCalibrationNet},
\texttt{FisheyeCornerLSTMSequenceCalibrationNet} oraz
\texttt{FisheyeCornerSSMSequenceCalibrationNet} --- zostały wytrenowane na
tych samych danych syntetycznych i ustawień, przy użyciu skryptów
\texttt{run\_fisheye\_corner\_sequence\_training.sh},
\texttt{run\_fisheye\_corner\_lstm\_sequence\_training.sh} oraz
\texttt{run\_fisheye\_corner\_ssm\_sequence\_training.sh}. Wszystkie trzy
uruchomienia treningu miały te same ustawienia: rozdzielczość
obrazu $1920\times1080$ px, plansza $10\times10$ pól (81 wewnętrznych
narożników), długość sekwencji $T=10$ klatek, zakresy kątów pochylenia
$\pm45^\circ$, odchylenia $\pm45^\circ$ i przechylenia $\pm30^\circ$,
optymalizator \texttt{AdamW} ze stopą uczenia $10^{-3}$, rozmiar paczki $64$,
maksymalnie $100$ epok oraz wczesne zatrzymywanie (\emph{early stopping}) przy braku poprawy wyników przez $10$ epok. Jedyną różnicą pomiędzy uruchomieniami był wybór
architektury (\texttt{--model-name}) oraz katalog wyjściowy checkpointu.

Ocenę modeli po treningu przeprowadzono w dwóch niezależnych krokach:
\begin{enumerate}[leftmargin=*]
  \item odczyt straty walidacyjnej i numeru najlepszej epoki zapisanych w
        punktach kontrolnych (\texttt{best\_model.pth}) każdego z trzech
        modeli,
  \item test na rzeczywistych zdjęciach z kamery, wykonany w notatnikach
        \texttt{detect\_chessboard\_corners.ipynb} oraz
        \texttt{compare\_fisheye\_calibrations.ipynb}, obejmujący detekcję
        narożników, klasyczną kalibrację referencyjną \texttt{cv2.fisheye} [5, 8]
        oraz predykcję parametrów przez każdy z trzech modeli na tej samej
        sekwencji rzeczywistych obrazów.
\end{enumerate}

\subsection{Wyniki treningu}
\label{sec:wyniki-treningu}

Tabela~\ref{tab:wyniki-treningu} zestawia epokę, w której zapisano najlepszy
checkpoint każdego modelu, oraz odpowiadającą jej wartość funkcji straty
\emph{log-cosh} na zbiorze walidacyjnym danych syntetycznych.

\begin{table}[htbp]
  \centering
  \begin{tabular}{@{}lccc@{}}
    \toprule
    \textbf{Model} & \textbf{Najlepsza epoka} & \textbf{Strata walid. (log-cosh)} & \textbf{Limit epok} \\
    \midrule
    GRU  & 28 & 0.000993 & 100 \\
    LSTM & 53 & \textbf{0.000978} & 100 \\
    SSM  & \textbf{12} & 0.001099 & 100 \\
    \bottomrule
  \end{tabular}
  \caption{Wyniki treningu trzech wariantów modelu sekwencyjnego (zbiór walidacyjny syntetyczny).}
  \label{tab:wyniki-treningu}
\end{table}

Model LSTM osiągnął najniższą stratę walidacyjną, jednak dopiero po $53$
epokach, co --- przy dwóch warstwach rekurencyjnych --- oznacza znacząco
dłuższy czas treningu niż w przypadku pozostałych wariantów. Model GRU
uzyskał zbliżoną, nieznacznie wyższą stratę po $28$ epokach. Model SSM
zatrzymał się najwcześniej (epoka $12$) z najwyższą spośród trzech wartością
straty walidacyjnej, co sugerowałoby na tym etapie najsłabszą jakość spośród
trzech wariantów. Jak pokazano w kolejnym podrozdziale, ranking oparty na
stracie walidacyjnej na danych syntetycznych \emph{nie} pokrywa się z
rankingiem uzyskanym na rzeczywistych zdjęciach z kamery.

\subsection{Wyniki na syntetycznych sekwencjach testowych}
\label{sec:wyniki-syntetyczne-sekwencje}

Niezależnie od testu na rzeczywistych zdjęciach przeprowadzono ewaluację na
dwóch zbiorach syntetycznych, z których każdy zawierał $20$ sekwencji po
$10$ klatek i $81$ narożników w każdej klatce. Pierwszy zbiór,
\texttt{fisheye\_corner\_test\_seq\_training\_range}, miał rozkład
parametrów zbliżony do zakresu wykorzystanego podczas treningu. Drugi zbiór,
\texttt{fisheye\_corner\_test\_seq}, obejmował silniejsze i szersze
zniekształcenia, dlatego był trudniejszym testem działania modelu poza zakresem treningu. W obu przypadkach
wartości referencyjne były znane, ponieważ punkty generowano bezpośrednio z
modelu projekcji fisheye.

W tabeli~\ref{tab:wyniki-syntetyczne-training-range} przedstawiono wyniki dla
zbioru o zakresie treningowym. Błąd parametrów jest średnim RMSE po ośmiu
parametrach, po ich normalizacji przez odpowiednio szerokość i wysokość
obrazu dla parametrów pikselowych. Błąd reprojekcji jest średnią odległością
punktu przewidzianego od punktu obserwowanego.

\begin{table}[htbp]
  \centering
  \small
  \begin{tabular}{@{}lccc@{}}
    \toprule
    \textbf{Metoda} & \textbf{RMSE parametrów norm.} & \textbf{MAE parametrów norm.} & \textbf{Błąd reprojekcji [px]} \\
    \midrule
    OpenCV fisheye & \textbf{0.000019} & \textbf{0.000011} & \textbf{0.000026} \\
    GRU            & 0.051401         & 0.040161         & 0.524564 \\
    LSTM           & 0.050684         & 0.039301         & \textbf{0.458812} \\
    SSM            & 0.050837         & 0.040681         & 0.498279 \\
    \bottomrule
  \end{tabular}
  \caption{Średnie wyniki kalibracji na 20 syntetycznych sekwencjach ze zbioru training range.}
  \label{tab:wyniki-syntetyczne-training-range}
\end{table}

W zakresie zbliżonym do danych treningowych klasyczna kalibracja OpenCV
odzyskała parametry niemal dokładnie, co jest oczekiwane dla bezszumowych
danych syntetycznych. Wśród modeli neuronowych najlepszy średni RMSE parametrów oraz błąd reprojekcji uzyskał LSTM. Wyniki tego eksperymentu pokazują, że implementacja działa poprawnie i że modele radzą sobie z wartościami pomiędzy przykładami treningowymi. Nie jest to jednak odpowiednik testu z rzeczywistymi danymi.

Na trudniejszym zbiorze \texttt{fisheye\_corner\_test\_seq} średnie RMSE
parametrów modeli wyniosło odpowiednio $0.125374$ dla GRU, $0.123439$ dla
LSTM oraz $0.118770$ dla SSM. W tym eksperymencie SSM uzyskał najlepszą
zgodność parametrów, co wskazuje na lepsze działanie poza podstawowym
zakres treningowy.

Klasyczna metoda OpenCV nie zachowywała się stabilnie na tym trudniejszym
zbiorze. Dla $9$ z $20$ sekwencji procedura inicjalizacji ekstrynsycznej
zgłaszała błąd \texttt{InitExtrinsics}. Po zastosowaniu awaryjnego pomijania
pojedynczej klatki część kalibracji kończyła się bez błędu, ale
zwracała nierealne parametry, na przykład bardzo duże współczynniki
zniekształcenia i błędne ogniskowe. Średni RMSE parametrów OpenCV wzrósł
wtedy do $50.751156$, przy maksimum $922.919015$. Oznacza to, że klasyczny
algorytm optymalizacyjny nie radził sobie z niektórymi kombinacjami silnej dystorsji i pozycji
planszy, a nie że punkty syntetyczne były niepoprawne. Poprawność generatora
potwierdzono przez ponowną projekcję punktów przy użyciu parametrów
referencyjnych, która dawała błąd rzędu $10^{-5}$~px.

Oprócz jakości zmierzono czas pojedynczego wywołania każdej metody. Pomiar
obejmował wyłącznie kalibrację lub inferencję dla jednej sekwencji, bez
ładowania checkpointu i inicjalizacji modelu. Wyniki przedstawiono w
tabeli~\ref{tab:czasy-syntetyczne}.

\begin{table}[htbp]
  \centering
  \small
  \begin{tabular}{@{}lcccc@{}}
    \toprule
    \textbf{Metoda} & \textbf{Średnia [s]} & \textbf{Odchylenie [s]} & \textbf{Minimum [s]} & \textbf{Maksimum [s]} \\
    \midrule
    OpenCV fisheye & 0.030350 & 0.010211 & 0.024726 & 0.063126 \\
    GRU            & 0.001284 & 0.002058 & 0.000692 & 0.010013 \\
    LSTM           & 0.000970 & 0.000296 & 0.000694 & 0.001785 \\
    SSM            & 0.001767 & 0.001350 & 0.001251 & 0.007464 \\
    \bottomrule
  \end{tabular}
  \caption{Czas pojedynczego wywołania metody na syntetycznej sekwencji, średnia z 20 sekwencji.}
  \label{tab:czasy-syntetyczne}
\end{table}

W badanym środowisku LSTM był najszybszym modelem neuronowym, natomiast
OpenCV wymagało średnio około $30$~ms na sekwencję. Odpowiada to około
$24$ razy dłuższemu czasowi niż dla GRU, około $31$ razy dłuższemu niż dla
LSTM i około $17$ razy dłuższemu niż dla SSM. Czasów tych nie należy jednak
interpretować jako uniwersalnego rankingu sprzętowego, ponieważ zależą od
procesora, wersji bibliotek i sposobu wywołania modelu. Pokazują one
natomiast różnicę między iteracyjną optymalizacją klasyczną a jednorazową
inferencją wytrenowanej sieci.

\subsection{Dane testowe z rzeczywistej kamery}
\label{sec:dane-realne}

Testy na rzeczywistej kamerze wykonano na $10$ zdjęciach szachownicy
$10\times10$ pól o boku $30$~mm, zarejestrowanych obiektywem szerokokątnym w
rozdzielczości $1920\times1080$ px. Algorytm \texttt{findChessboardCorners}
z doprecyzowaniem \texttt{cornerSubPix} wykrył komplet $81$ narożników na
każdym zdjęciu. Walidacja końcowa potwierdziła brak nieudanych detekcji.

Na tak przygotowanych punktach wykonano referencyjną kalibrację
\texttt{cv2.fisheye.calibrate} [8], uzyskując globalny błąd reprojekcji
$\mathrm{RMSE} = 0.4232$~px (dla poszczególnych zdjęć od $0.344$~px do
$0.683$~px), przy macierzy kamery
\begin{equation}
  K_{\text{OpenCV}} = \begin{bmatrix}
    1046.69 & 0 & 958.08 \\
    0 & 1047.87 & 500.37 \\
    0 & 0 & 1
  \end{bmatrix}
\end{equation}
oraz współczynnikach dystorsji $[k_1,k_2,k_3,k_4] = [-0.06285,\ 0.00334,\
-0.01334,\ 0.00473]$. Wynik ten stanowi punkt odniesienia
dla oceny predykcji sieci neuronowych w kolejnym podrozdziale.

\subsection{Wyniki predykcji sieci na rzeczywistej kamerze}
\label{sec:wyniki-realne}

Ta sama sekwencja $10$ rzeczywistych zdjęć (w postaci znormalizowanych i
ustandaryzowanych cech narożników) została podana na wejście
każdego z trzech wytrenowanych modeli. Tabela~\ref{tab:wyniki-parametry}
zestawia parametry kalibracji wraz z błędem reprojekcji, przy
czym błąd reprojekcji sieci obliczono metodą spójną dla wszystkich wariantów:
punkty obrazu poddano \texttt{cv2.fisheye.undistortPoints} z użyciem
przewidzianych $K$ i $D$, wyznaczono pozę planszy funkcją
\texttt{cv2.solvePnP} (kamera jednostkowa, brak dystorsji), a następnie
rzutowano punkty 3D z powrotem na obraz funkcją
\texttt{cv2.fisheye.projectPoints} [8] przy użyciu przewidzianych $K$ i $D$.

\begin{table}[htbp]
  \centering
  \small
  \begin{tabular}{@{}lccccc@{}}
    \toprule
    \textbf{Parametr} & \textbf{OpenCV fisheye} & \textbf{GRU} & \textbf{LSTM} & \textbf{SSM} \\
    \midrule
    $f_x$ [px] & 1046.69 & 1050.54 & 1091.33 & 1042.73 \\
    $f_y$ [px] & 1047.87 & 1130.83 & 1116.75 & 1041.87 \\
    $c_x$ [px] & 958.08  & 953.35  & 956.05  & 961.52  \\
    $c_y$ [px] & 500.37  & 490.28  & 495.73  & 499.63  \\
    $k_1$      & $-0.06285$ & $-0.00460$ & $-0.01237$ & $-0.06070$ \\
    $k_2$      & $0.00334$  & $0.02837$  & $0.04254$  & $0.00443$  \\
    $k_3$      & $-0.01334$ & $-0.00109$ & $0.01457$  & $-0.01246$ \\
    $k_4$      & $0.00473$  & $0.01751$  & $0.02388$  & $0.00347$  \\
    \midrule
    \textbf{RMSE reprojekcji [px]} & \textbf{0.4232} & 6.6312 & 3.0981 & \textbf{0.4799} \\
    \bottomrule
  \end{tabular}
  \caption{Parametry kalibracji przewidziane przez sieci GRU, LSTM i SSM na 10 rzeczywistych zdjęciach, na tle referencji OpenCV fisheye.}
  \label{tab:wyniki-parametry}
\end{table}

Tabela~\ref{tab:wyniki-roznice} przedstawia te same wyniki w formie różnic
bezwzględnych każdego modelu względem kalibracji referencyjnej OpenCV, co
ułatwia ocenę, który wariant najwierniej odtwarza rzeczywiste parametry
kamery.

\begin{table}[htbp]
  \centering
  \small
  \begin{tabular}{@{}lccc@{}}
    \toprule
    \textbf{Parametr} & \textbf{$|$GRU $-$ OpenCV$|$} & \textbf{$|$LSTM $-$ OpenCV$|$} & \textbf{$|$SSM $-$ OpenCV$|$} \\
    \midrule
    $f_x$ [px] & \textbf{3.85}  & 44.64 & 3.96 \\
    $f_y$ [px] & 82.96 & 68.87 & \textbf{6.00} \\
    $c_x$ [px] & 4.73  & \textbf{2.03}  & 3.44 \\
    $c_y$ [px] & 10.09 & 4.65  & \textbf{0.74} \\
    $k_1$      & 0.05825 & 0.05048 & \textbf{0.00215} \\
    $k_2$      & 0.02503 & 0.03920 & \textbf{0.00109} \\
    $k_3$      & 0.01225 & 0.02791 & \textbf{0.00089} \\
    $k_4$      & 0.01278 & 0.01916 & \textbf{0.00126} \\
    \bottomrule
  \end{tabular}
  \caption{Różnice bezwzględne parametrów sieci względem kalibracji OpenCV fisheye (rzeczywista kamera).}
  \label{tab:wyniki-roznice}

  \vspace{2pt}
\end{table}

Wyniki na rzeczywistej kamerze wskazują jednoznacznie, że model SSM ---
mimo najniższej pojemności (najmniej parametrów, patrz
tabela~\ref{tab:gru-lstm-ssm}) i najwyższej straty walidacyjnej na danych
syntetycznych --- najlepiej przenosi się na rzeczywiste zdjęcia spośród
trzech wariantów, osiągając błąd reprojekcji ($0.48$~px) niemal identyczny z
klasyczną kalibracją OpenCV ($0.42$~px) oraz najmniejsze odchyłki dla
$f_y$, $c_y$ i wszystkich współczynników dystorsji. Model GRU i model LSTM
przewidują natomiast wyraźnie zawyżoną ogniskową $f_y$ (odpowiednio o
$83$~px i $69$~px) oraz współczynniki dystorsji radialnej bliskie zeru
($k_1 \approx -0.005$ i $-0.012$ zamiast $-0.063$ referencyjnego), co
przekłada się na błędy reprojekcji wielokrotnie większe niż w przypadku
SSM i OpenCV.

Jako wizualne uzupełnienie porównania wykonano także undystorsję pojedynczego
zdjęcia \texttt{capture\_0200\_correct.png} przy użyciu parametrów
uzyskanych przez wszystkie cztery metody. Na rysunkach
\ref{fig:undystorsja-opencv}--\ref{fig:undystorsja-ssm} po lewej stronie
pokazano obraz wejściowy, a po prawej obraz po korekcji dystorsji. Zestawienie
pozwala ocenić nie tylko wartości liczbowe błędu reprojekcji, lecz także
wizualny wpływ przewidzianych parametrów na geometrię obrazu.

\begin{figure}[htbp]
  \centering
  \includegraphics[width=0.98\textwidth]{img/capture_0200_OpenCV_fisheye.png}
  \caption{Para obrazu przed undystorsją i po undystorsji uzyskanej z parametrów kalibracji OpenCV fisheye.}
  \label{fig:undystorsja-opencv}
\end{figure}

\begin{figure}[htbp]
  \centering
  \includegraphics[width=0.98\textwidth]{img/capture_0200_Siec_GRU.png}
  \caption{Para obrazu przed undystorsją i po undystorsji uzyskanej z parametrów modelu GRU.}
  \label{fig:undystorsja-gru}
\end{figure}

\begin{figure}[htbp]
  \centering
  \includegraphics[width=0.98\textwidth]{img/capture_0200_Siec_LSTM.png}
  \caption{Para obrazu przed undystorsją i po undystorsji uzyskanej z parametrów modelu LSTM.}
  \label{fig:undystorsja-lstm}
\end{figure}

\begin{figure}[htbp]
  \centering
  \includegraphics[width=0.98\textwidth]{img/capture_0200_Siec_SSM.png}
  \caption{Para obrazu przed undystorsją i po undystorsji uzyskanej z parametrów modelu SSM.}
  \label{fig:undystorsja-ssm}
\end{figure}

\subsection{Uwaga metodologiczna dotycząca porównania z klasyczną kalibracją}
\label{sec:uwaga-metodologiczna}

Podczas przygotowywania niniejszego rozdziału wykryto i poprawiono usterkę w
funkcji \texttt{rmse} notatnika \texttt{compare\_fisheye\_calibrations.ipynb}
(oraz odpowiadającego mu skryptu
\texttt{compare\_fisheye\_calibrations.py}). W pierwotnej wersji funkcja ta,
w gałęzi obsługującej wiersz \texttt{OpenCV fisheye}, przekazywała
współczynniki dystorsji rybiego oka $[k_1,k_2,k_3,k_4]$ bezpośrednio do
\texttt{cv2.solvePnP} jako parametr \texttt{distCoeffs}. Funkcja
\texttt{cv2.solvePnP} interpretuje jednak ten parametr zgodnie z konwencją
modelu Browna--Conrady'ego ($k_1,k_2,p_1,p_2[,k_3,\dots]$), a nie modelu
rybiego oka, co prowadziło do błędnego wyznaczenia pozy planszy i sztucznie
zawyżonego błędu reprojekcji rzędu $44$~px dla tego jednego wiersza tabeli
porównawczej --- w oczywistej sprzeczności z wartością $0.4232$~px
wyznaczoną w tym samym projekcie, w notatniku
\texttt{detect\_chessboard\_corners.ipynb}, przy użyciu wektorów rotacji i
translacji zwróconych bezpośrednio przez \texttt{cv2.fisheye.calibrate}.

Usterkę usunięto, stosując tę samą procedurę dla wszystkich czterech metod:
punkty obrazu są zawsze najpierw przekształcane funkcją
\texttt{cv2.fisheye.undistortPoints} przy użyciu macierzy $K$ i
współczynników $D$ danej metody, a dopiero tak znormalizowane punkty trafiają
do \texttt{cv2.solvePnP} wywoływanego z kamerą jednostkową i zerową
dystorsją --- identycznie jak było to już wcześniej realizowane dla modeli
GRU, LSTM i SSM. Po poprawce funkcja rmse zwraca dla metody OpenCV
fisheye wartość $0.4235$~px, praktycznie identyczną z referencyjną wartością
$0.4232$~px z notatnika kalibracyjnego (różnica $0.0003$~px wynika z tego, że
w \texttt{compare\_fisheye\_calibrations.ipynb} poza planszy jest
każdorazowo wyznaczana na nowo funkcją \texttt{solvePnP}, zamiast wykorzystywać
wektory rotacji i translacji zwrócone bezpośrednio przez
\texttt{cv2.fisheye.calibrate}). Potwierdza to, że wartości referencyjne
OpenCV użyte w tabelach~\ref{tab:wyniki-parametry} i~\ref{tab:wyniki-roznice}
niniejszego rozdziału ($0.4232$~px) są poprawne, a wnioski z
podrozdziału~\ref{sec:wyniki-realne} --- w szczególności bliskość wyniku
modelu SSM ($0.4799$~px) do wyniku klasycznej kalibracji --- pozostają w
mocy po naprawieniu błędu w kodzie porównawczym.

% koniec pliku

